\chapter*{Acknowledgements}
\addcontentsline{toc}{chapter}{\protect\numberline{}Acknowledgements}

Over the years, I have been asked a number of times whether I would recommend anyone to do a PhD. My answer has always been a resounding "No.", followed by a brief pause and "...unless you can find a really good lab". What makes a good lab? First of all, there must be infrastructure and funding available to support research. This is easy for a prospective student to verify: go and visit, and take a look. Second, the staff and students should be friendly and intelligent. You will draw on them for inspiration and companionship, and the effects of good colleagues on your own work will be profound. Finally, the person in charge should be a visionary. Ideally, he or she has clear and unique ideas of what the future will look like, and how the lab can contribute to and interact with it. Only then novel output can be produced consistently.

I have been fortunate to have ended up, mostly by chance, in a lab where all three conditions were exceedingly met. At Aerts lab, I have met some of the kindest and most intelligent people, who made research fun and inspiring. During (and after) my tenure there, funding was ample, and we were able to invest time and money into crazy, low-percentage ideas - which eventually did work out. Admittedly, most of the time we were just finding and doing things for the first time ourselves. All of this is downstream of good leadership, with the courage to commit fully to unique ideas - effectively betting the entire lab's existence on a small number of powerful hypotheses.

I would like to thank my advisor Stein, for building this environment. I am lucky and grateful to be able to write this acknowledgement two years after leaving the lab, with newfound appreciation for the environment in which my career was started. I would also like to thank my supervisor Suresh for his excellent daily support. We endured a lot together in building HyDrop! Finally, I would also like to thank Jasper for his supervision during my master thesis and introducing me to the lab.